\section{Ordinary reduction and explicit limits of the filtration condition}\label{sec:early-reduction}
This section proves the ordinary good-reduction criterion, gives actual elliptic examples with invisible first-level special characters, and explains why good graded pieces do not force potential semistability of the entire representation.

\subsection{Ordinary automorphisms and unramified torsion}
The ordinary criterion rests on two elementary faithfulness arguments.

\begin{lemma}\label{lem:principal-congruence-torsion}
For odd $p$, the kernel of $\GL_g(\Z_p)\to\GL_g(\F_p)$ is torsion-free.
\end{lemma}
\begin{proof}
Write a nonidentity element as $U=1+p^aB$, where $a\ge1$ and $B$ is nonzero modulo $p$. If $r$ is prime to $p$, the binomial expansion gives
\[
 U^r=1+p^a(rB+pC),\qquad
 U^p=1+p^{a+1}(B+pC')
\]
with integral matrices $C,C'$. In the second formula, the middle binomial terms have valuation at least $a+2$, and the last term has valuation $ap\ge a+2$. Thus taking a prime-to-$p$ power preserves the least matrix valuation, while taking a $p$th power raises it by exactly one. An element of finite order $p^br$ would become the identity after these operations, a contradiction.
\end{proof}

\begin{lemma}\label{lem:ordinary-etale-faithful}
Let $B$ be an ordinary abelian variety of dimension $g$ over an algebraically closed field $k$ of characteristic $p$. An endomorphism annihilating $B[p^m](k)$ for every $m$ is zero. Consequently $\operatorname{Aut}(B,0)$ acts faithfully on $T_p^{\mathrm{et}}B$.
\end{lemma}
\begin{proof}
For an abelian variety $C/k$ of dimension $h$, the factorization $[p^m]=V^mF^m$ has $F^m$ radicial and bijective on geometric points, and $V^m$ of degree $p^{mh}$. Hence
\begin{equation}\label{eq:ordinary-point-count}
 |C[p^m](k)|\le p^{mh},\qquad |B[p^m](k)|=p^{mg}.
\end{equation}
The equality is the ordinary condition. The factorization and degree are the standard Frobenius--Verschiebung statements of \citet[Proposition~5.20, paragraph~5.21 and Proposition~5.22, pp.~80--81]{EdixhovenGeerMoonen}.

Suppose $f\ne0$ annihilates all these points. The reduced identity component $H$ of $\ker(f)$ is an abelian subvariety of dimension $h<g$. Indeed the reduced subgroup over the perfect field $k$ is a group scheme; translation carries its nonempty smooth locus to every point, so it is smooth. Its identity component is connected and proper. If it had dimension $g$, then $f$ would vanish on all of $B$. The reduced kernel has finitely many components, say $c$, each a translate of $H$.

Each such translate contains at most $|H[p^m](k)|$ points of $B[p^m](k)$: if it contains one, subtracting that point identifies the intersection with $H[p^m](k)$. Equation~\eqref{eq:ordinary-point-count} would therefore give $p^{mg}\le c p^{mh}$ for all $m$, impossible for $h<g$. Thus $f=0$. If an automorphism is the identity on $T_p^{\mathrm{et}}B$, apply this assertion to its difference from the identity. The finite quotients of that Tate module are precisely $B[p^m](k)$.
\end{proof}

\begin{proof}[Proof of Theorem~\ref{thm:ordinary-unramified-p}]
The dimension-zero case is immediate. Pass to the strict henselian unramified base to describe inertia, and choose a finite Galois extension $L$ over which $A$ has good ordinary reduction, with abelian model $\mathcal A$ and special fibre $B$. Generic descent isomorphisms extend uniquely to the good models by the N\'eron mapping property \citep[Theorem~17.1 and Remark~17.2, p.~70]{Milne2008AV}. Their cocycle identities extend by uniqueness. Since inertia fixes the residue field, reduction gives a finite group of origin-preserving automorphisms of $B$.

The ordinary connected--\'etale sequence of $\mathcal A[p^\infty]$ is functorial under these isomorphisms. Its \'etale quotient has height $g=\dim A$, and its generic $p$-torsion is an equivariant quotient of $A[p]$. The unramified hypothesis makes inertia trivial on that quotient. Finite \'etale specialization over the henselian base identifies it with $B[p](k)=T_p^{\mathrm{et}}B/p$. Thus the finite special inertia image in $\GL_g(\Z_p)$ is congruent to the identity modulo $p$.

Lemma~\ref{lem:principal-congruence-torsion} makes this image trivial; Lemma~\ref{lem:ordinary-etale-faithful} then makes every special inertia automorphism of $B$ trivial. Smooth proper base change identifies the action on $T_\ell B$ with the original inertia action on $T_\ell A$, for $\ell\ne p$. The N\'eron--Ogg--Shafarevich criterion over the original $K$ now gives good reduction \citep[Theorem~1, p.~493]{SerreTate1968}. No cyclotomic determinant or polarization was used, so no roots-of-unity or polarization-degree hypothesis is needed.
\end{proof}

\subsection{A tame defect-two example at every odd prime}
The first family is a direct specialization of the complete tame existence theorem.

\begin{proof}[Proof of Corollary~\ref{cor:defect-two-full-p}]
In Theorem~\ref{thm:tame-level-one} take $m=2$ and $e=(p^2-1)/2$. For odd $p$, one has $2\mid p+1$ and $(p-1)\mid e$, and the lower endpoint is attained. The tame $p=3$ clause gives the same conclusion at $e=4$. The construction in Section~\ref{sec:tame-full-torsion} gives supersingular reduction; alternatively Theorem~\ref{thm:ordinary-unramified-p} excludes the ordinary case for an additive curve with full rational $p$-torsion.

If $P,Q$ form a basis of $E[p](K)$, the perfect alternating Weil pairing $e_p(P,Q)$ has order $p$ and is $K$-rational by equivariance \citep[Section~I.13, p.~57]{Milne2008AV}. Thus $\mu_p\subset K$ without enlarging the field. The nontrivial special inertia automorphism $\alpha$ has order two and fixes the origin. Since $(\alpha-1)(\alpha+1)=0$ and a nonzero elliptic endomorphism is a finite surjective isogeny, $\alpha\ne1$ forces $\alpha=-1$. Its trace on the rank-two rational Dieudonn\'e module is $-2$.
\end{proof}

\subsection{An explicit defect-six model in characteristic three}
The next calculation determines both the full torsion and the exact good-reduction field of one curve.

\begin{proof}[Proof of Proposition~\ref{prop:explicit-cubic-defect-six}]
Put $\zeta=\zeta_3$, $\tau=\zeta-\zeta^2=2\zeta+1$. Then $\tau^2=-3$, $v_K(\tau)=1$ and $v_K(3)=2$. The standard invariants of $y^2+y=x^3$ give $\Delta=-27$ and
\[
                      \psi_3(x)=3x^4+3x.
\]
Its eight nonidentity $3$-torsion points are
\[
 (0,0),\ (0,-1),\quad
 (-\zeta^i,\zeta),\ (-\zeta^i,\zeta^2)\quad(0\le i\le2).
\]
They satisfy the equation and the division polynomial, are distinct, and together with the identity exhaust the nine geometric $3$-torsion points in characteristic zero.

The displayed model is minimal because $v_K(\Delta)=6$ and a change of Weierstrass model changes that valuation by an integer multiple of twelve. A smaller nonnegative value is impossible. The reduced singular point is $(-1,1)$; the substitution $X=x+1$, $Y=y-1$ gives $Y^2=X^3$. This is a cusp and has additive smooth group. Explicitly its smooth locus has parameter $w=X/Y$, with $w=0$ at infinity and $(X,Y)=(w^{-2},w^{-3})$ away from infinity. A line gives an equation $cw^3+aw+b=0$, whose three roots sum to zero, so the chord law is addition of $w$.

Choose $r^3=-1/4$ and $u^2=\tau$, and put $L=K(r,u)$. The unit $r$ reduces to $-1$. If $s=r+1$, then
\[
                     s^3-3s^2+3s-3/4=0.
\]
For $t=v_K(s)>0$, the four term valuations are $3t,2+2t,2+t,2$. A vanishing sum cannot have a unique term of least valuation, so $3t=2$. Thus $K(r)/K$ has ramification index and degree three; it is cyclic because $\mu_3\subset K$. Also $v_K(u)=1/2$, giving a ramified quadratic extension. Their compositum is cyclic of degree six and has ramification index six, since its index is divisible by both two and three.

Under
\[
                    x=r+\tau X,\qquad y=u^3Y-1/2,
\]
the equation becomes
\begin{equation}\label{eq:explicit-cubic-good-model}
                      Y^2=X^3-\tau rX^2-r^2X.
\end{equation}
Its coefficients are integral and its discriminant is
\[
 16r^4(\tau^2r^2+4r^2)=16r^6(\tau^2+4)=16r^6=1.
\]
This is an actual good model. Its special fibre is $E_0:Y^2=X^3-X$. With the standard convention
\[
 \psi_3=3X^4+b_2X^3+3b_4X^2+3b_6X+b_8,
\]
one has $b_8=-1$, and hence $\psi_3=-1$ in characteristic three. There is no nonidentity geometric $3$-torsion point; the elliptic curve has $p$-rank zero and is supersingular. Equivalently, for $Y\ne0$ the tangent slope is $-1/(2Y)$, and the doubling formula gives $X(2P)=X(P)+1/(4Y(P)^2)\ne X(P)$, excluding $2P=-P$; points with $Y=0$ have order two.

The commuting generators $\sigma_3(r)=\zeta r$, $\sigma_3(u)=u$ and $\sigma_2(r)=r$, $\sigma_2(u)=-u$ act in good coordinates by
\[
 \sigma_3(X)=X+\frac{(1-\zeta)r}{\tau}=X+\zeta^{-1}r,
 \quad \sigma_3(Y)=Y,
 \qquad \sigma_2(X)=X,\quad\sigma_2(Y)=-Y.
\]
After reduction these are $(X,Y)\mapsto(X-1,Y)$ and $(X,Y)\mapsto(X,-Y)=[-1]$. The first cyclically permutes the three nonzero $2$-torsion points $(0,0),(1,0),(-1,0)$, so is nontrivial on $T_2E_0$. The second is $-1$ on that Tate module. Any nontrivial subgroup of the cyclic group of order six contains its subgroup of order two or three, so this action is faithful.

Smooth proper base change therefore makes the inertia image on $T_2E$ have exact order six. If $E$ becomes good over $F/K$, NOS makes $I_F$ lie in the kernel of that quotient, so $6\mid[I_K:I_F]=e(F/K)$. The field $L$ attains six. Since $j(E)=0$, multiplicative reduction is impossible after any extension: a nodal minimal cubic has $c_4$ a unit and $v(j)=3v(c_4)-v(\Delta)<0$. Hence good reduction and semistability have the same minimal ramification here.

After unramified base change, $v(\tau)=1$ and $v(3)=2$ remain unchanged. The same two valuation calculations and the same faithful special action still apply. This proves the exact defect also over the maximal unramified base and its completion. The special involution is $[-1]$ and has trace $-2$.
\end{proof}

\subsection{The logarithmic cyclotomic extension}
The final example explains the representation-theoretic scope of the allowed filtration class.

\begin{proof}[Proof of Proposition~\ref{prop:filtered-non-ht}]
Let $\chi_p:G_K\to\Z_p^\times$ be the cyclotomic character. Since $K=\Q_p(\zeta_p)$, its values lie in $1+p\Z_p$. The convergent logarithm gives the additive continuous function $c(g)=\log_p\chi_p(g)\in p\Z_p$. Set
\begin{equation}\label{eq:log-cyclotomic-representation}
 T=\Z_pe_1\oplus\Z_pe_2,\qquad
 \rho(g)=\begin{pmatrix}1&c(g)\\0&1\end{pmatrix}.
\end{equation}
Additivity of $c$ proves the representation identity. The submodule $\Z_pe_1$ and its quotient are trivial Galois modules.

Let $G_m$ be the finite \'etale $K$-group scheme corresponding to $T/p^mT$. The equivariant injections $a\mapsto pa$ from $T/p^mT$ to $T/p^{m+1}T$ give a Barsotti--Tate system of height two. Indeed the kernel of $p^r$ on $T/p^{m+r}T$ is $p^mT/p^{m+r}T$, the image of $T/p^rT$ under the indicated transition. The invariant line and quotient at every level are constant and compatible, giving the exact sequence in the proposition. Its Tate module is $T$.

Put $N(e_2)=e_1$, $N(e_1)=0$. Then $N\ne0$, $N^2=0$ and $\rho(g)=\exp(c(g)N)$. The cyclotomic kernel acts trivially. In the cyclotomic Sen module, choose a sufficiently small $\gamma$ fixing the finite coefficient field of the basis. The normalized logarithm is
\[
               \Theta=\frac{\log\rho(\gamma)}{\log_p\chi_p(\gamma)}=N.
\]
The operator is nonsemisimple, so Sen's criterion excludes Hodge--Tateness. Its compatibility with finite extension leaves the same nonzero nilpotent operator after every finite base extension, excluding potential Hodge--Tateness. Since semistable representations are de Rham and de Rham representations are Hodge--Tate, potential semistability is also excluded.

This is precisely the classical example in \citet[Section~II.1.2, p.~9 of the 2002 author version]{Berger2004}. Each graded piece of the constructed group extends to the constant height-one \'etale group over $\Ocal_K$ and has special inertia character one. Their sum is the constant character two. This graded character does not imply a potentially good realization of the whole group.
\end{proof}
